% Figure: Immune Cell Dynamics and Cytokine Network Models
% (a) NK cell and T cell state transition diagrams
% (b) Cytokine network topology

\begin{figure}[htbp]
\centering
\begin{tikzpicture}[
    node distance=2.5cm and 3cm,
    scale=0.82, every node/.style={scale=0.82},
    % Cell state styles
    resting/.style={draw=green!70!black, fill=green!10, very thick, rounded corners,
        minimum width=2.2cm, minimum height=1cm, align=center},
    active/.style={draw=blue!70!black, fill=blue!10, very thick, rounded corners,
        minimum width=2.2cm, minimum height=1cm, align=center},
    exhausted/.style={draw=red!70!black, fill=red!10, very thick, rounded corners,
        minimum width=2.2cm, minimum height=1cm, align=center},
    memory/.style={draw=teal!70!black, fill=teal!10, very thick, rounded corners,
        minimum width=2.2cm, minimum height=1cm, align=center},
    % Arrow styles
    trans/.style={-latex, very thick, line width=1.2pt},
    slow/.style={-latex, very thick, line width=1.2pt, dashed},
    death/.style={-latex, thin, gray!70, line width=0.8pt},
    input/.style={-latex, thick, green!60!black, line width=1pt},
    % Cytokine node styles
    proinflam/.style={circle, draw=red!70!black, fill=red!10, very thick,
        minimum size=1.3cm, align=center, font=\small\bfseries},
    antiinflam/.style={circle, draw=blue!70!black, fill=blue!10, very thick,
        minimum size=1.3cm, align=center, font=\small\bfseries},
    % Cytokine arrow styles
    stim/.style={-latex, thick, green!60!black, line width=1pt},
    inhib/.style={-|, thick, red!70!black, line width=1pt},
    % Labels
    sublabel/.style={font=\large\bfseries},
    ratelabel/.style={font=\scriptsize, fill=white, inner sep=1.5pt},
]

% ============================================================
% PART (a): Cell State Transition Diagrams
% ============================================================

\node[sublabel] at (-4.5, 0.8) {(a)};

% --- NK Cell Model (left side) ---
\node[font=\bfseries] at (-4.5, 0) {NK Cell Model (Eq.~28.1)};

\node[resting] (Nr) at (-7, -1.8) {$N_r$\\[-2pt]{\scriptsize Resting}};
\node[active]  (Na) at (-2, -1.8) {$N_a$\\[-2pt]{\scriptsize Activated}};
\node[exhausted] (Ne) at (-4.5, -4.5) {$N_e$\\[-2pt]{\scriptsize Exhausted}};

% Transitions
\draw[trans] (Nr) -- node[ratelabel, above] {$k_{\text{act}}(C)$} (Na);
\draw[trans] (Na) -- node[ratelabel, right=1pt] {$k_{\text{exh}}$} (Ne);
\draw[slow]  (Ne) -- node[ratelabel, left=1pt] {$k_{\text{recov}}$} (Nr);

% Input from bone marrow
\draw[input] (-7, 0) -- node[ratelabel, right=1pt] {$s_N$} (Nr);
\node[font=\scriptsize\itshape, text=green!50!black] at (-7, 0.3) {bone marrow};

% Death arrows
\draw[death] (Nr) -- ++(-0.9, -0.7) node[font=\tiny, gray, below] {$d_r$};
\draw[death] (Na) -- ++(0.9, -0.7) node[font=\tiny, gray, below] {$d_a$};
\draw[death] (Ne) -- ++(0, -1.0) node[font=\tiny, gray, below] {$d_e$};

% --- T Cell Model (right side) ---
\node[font=\bfseries] at (4.5, 0) {T Cell Model (Eq.~28.4)};

\node[resting] (Tn) at (1.5, -1.8) {$T_n$\\[-2pt]{\scriptsize Na\"ive}};
\node[active]  (Te) at (5.5, -1.8) {$T_e$\\[-2pt]{\scriptsize Effector}};
\node[memory]  (Tm) at (7.5, -4.5) {$T_m$\\[-2pt]{\scriptsize Memory}};
\node[exhausted] (Tex) at (3.5, -4.5) {$T_{ex}$\\[-2pt]{\scriptsize Exhausted}};

% Transitions
\draw[trans] (Tn) -- node[ratelabel, above] {$k_{\text{prime}} \!\cdot\! A(t)$} (Te);
\draw[trans] (Te) -- node[ratelabel, right=1pt] {$k_{\text{mem}}$} (Tm);
\draw[trans] (Te) -- node[ratelabel, left=1pt] {$k_{T_{ex}}$} (Tex);
\draw[slow]  (Tm) to[bend right=25] node[ratelabel, below] {\scriptsize reactivation} (Te);

% Self-loop for proliferation
\draw[trans] (Te) to[out=60, in=120, looseness=4]
    node[ratelabel, above=2pt] {$\rho_T$} (Te);

% ============================================================
% PART (b): Cytokine Network Topology
% ============================================================

\node[sublabel] at (-7, -6.5) {(b)};
\node[font=\bfseries] at (-1, -6.5) {Cytokine Network Topology};

% Pro-inflammatory nodes (warm colors)
\node[proinflam] (tnfa) at (-5.5, -8.8) {TNF-$\alpha$};
\node[proinflam] (il1b) at (-2.5, -8.8) {IL-1$\beta$};
\node[proinflam, draw=orange!70!black, fill=orange!10] (il6) at (0.5, -8.8) {IL-6};
\node[proinflam, draw=red!60!black, fill=red!8] (ifng) at (3.5, -8.8) {IFN-$\gamma$};

% Anti-inflammatory nodes (cool colors)
\node[antiinflam] (il10) at (-3.5, -11.8) {IL-10};
\node[antiinflam, draw=teal!70!black, fill=teal!10] (tgfb) at (1.5, -11.8) {TGF-$\beta$};

% Source labels
\node[font=\scriptsize\itshape, text=gray!70!black, above=0.4cm of tnfa] {$M_a$ (macrophages)};
\node[font=\scriptsize\itshape, text=gray!70!black, above=0.4cm of ifng] {T cells};

% Stimulatory arrows
\draw[stim] (tnfa) -- node[ratelabel, above, font=\tiny] {stim.} (il6);
\draw[stim] (il1b) -- node[ratelabel, above, font=\tiny] {stim.} (il6);

% IL-6 positive feedback loop via Th17
\draw[stim] (il6) to[out=-30, in=30, looseness=3.5]
    node[ratelabel, right=2pt, font=\tiny, text width=1cm, align=left] {Th17\\feedback} (il6);

% IFN-γ → IDO note
\draw[stim] (ifng) -- ++(1.8, 0)
    node[ratelabel, right=-1pt, font=\tiny] {$\to$ IDO};

% Inhibitory arrows (flat-headed, red)
\draw[inhib] (il10) to[bend right=12] (tnfa);
\draw[inhib] (il10) -- (il1b);
\draw[inhib] (il10) to[bend left=12] (il6);
\draw[inhib] (tgfb) to[bend left=15] (ifng);

% Legend
\begin{scope}[shift={(-7, -13.5)}]
    \draw[stim] (0,0) -- (0.8,0) node[right, font=\scriptsize] {Stimulation};
    \draw[inhib] (3,0) -- (3.8,0) node[right, font=\scriptsize] {Inhibition};
    \draw[slow] (6,0) -- (6.8,0) node[right, font=\scriptsize] {Slow / impaired in ME/CFS};
\end{scope}

\end{tikzpicture}
\caption{Immune system model components. (a)~State transition diagrams for NK cells (left, Equation~\ref{eq:nk-dynamics}) and T~cells (right, Equation~\ref{eq:tcell-dynamics}). In ME/CFS, accelerated exhaustion ($k_{\text{exh}}$ increased) and impaired recovery ($k_{\text{recov}}$ reduced) shift populations toward exhausted states (red). (b)~Cytokine network topology showing pro-inflammatory (warm colors) and anti-inflammatory (cool colors) mediators with stimulatory (arrows) and inhibitory (flat-headed) interactions.}
\label{fig:immune-cell-dynamics}
\end{figure}
